\documentclass[11pt]{article}

\usepackage[margin=1in]{geometry}
\usepackage{amsmath,amssymb}
\usepackage{bm}
\usepackage[hidelinks]{hyperref}

\newcommand{\kv}{\mathbf{k}}
\newcommand{\qv}{\mathbf{q}}
\newcommand{\Gv}{\mathbf{G}}
\newcommand{\Rv}{\mathbf{R}}
\newcommand{\Zst}{\mathbf{Z}^{*}}
\newcommand{\epsinf}{\bm{\varepsilon}_{\infty}}
\newcommand{\gshort}{g^{\mathcal{S}}}
\newcommand{\glong}{g^{\mathcal{L}}}
\newcommand{\Dlong}{D^{\mathcal{L}}}
\newcommand{\Qv}{\mathbf{Q}}

\title{An \emph{ab-initio} Diagrammatic Monte Carlo prototype\\
       for the LiF electron polaron}
\author{Samuel Recker}
\date{\today}

\begin{document}
\maketitle

\begin{abstract}
This note fixes the \textbf{goals}, deliberate \textbf{non-goals}, and
\textbf{deferred} extensions of the \texttt{rmc-ab-initio} project: a diagrammatic
Monte Carlo (DiagMC) prototype, built on the \texttt{ironmc} engine, that replaces the
analytic Fr\"ohlich electron--phonon vertex $|g(\qv)|^{2}\propto\alpha/q^{2}$ with the
\emph{real} electron--phonon matrix elements $g(\kv,\qv,\nu)$ obtained from
density-functional perturbation theory and Wannier interpolation (Quantum \textsc{espresso}
$+$ \textsc{epw}). The first validation material is the electron polaron in LiF, a
wide-gap ionic insulator that sits close to the textbook Fr\"ohlich limit and therefore
provides a built-in analytic cross-check. The scope below is the \emph{first} prototype;
Section~\ref{sec:deferred} records what is intentionally postponed.
\end{abstract}

\section{Overview}
\label{sec:overview}

\texttt{ironmc} is a Rust Monte Carlo framework; its DiagMC crates sample the imaginary-time
Green's function of a polaron. This project (the \texttt{ironmc-ab-initio} repository, built
on the framework) uses \texttt{rmc-diagmc}, a state-generic DiagMC engine exposing a \texttt{Model} trait (electron dispersion, phonon dispersion, complex e--ph
vertex, phonon-momentum proposal), plus a standalone \emph{ab-initio} consumer
\texttt{rmc-ab-initio}. Two models implement the trait: the closed-form \emph{analytic}
Fr\"ohlich model, and the \emph{ab-initio} model whose material-specific electron--phonon
(e--ph) vertex $g(\kv,\qv,\nu)$ is read from a real DFT/DFPT + Wannier calculation. The
ground-state energy is read off the large-$\tau$ decay of $G(\tau)$; there is no
self-consistent iteration.

The data pipeline is
\[
  \underbrace{\text{pw.x},\ \text{ph.x}}_{\text{Quantum \textsc{espresso}}}
  \;\longrightarrow\;
  \underbrace{\text{epw.x}}_{\text{Wannier e--ph tensors}}
  \;\longrightarrow\;
  \underbrace{\text{Rust reader + reconstruction}}_{\texttt{ironmc-ab-initio}}
  \;\longrightarrow\;
  \text{DiagMC}.
\]
DFPT supplies the phonons, the Born effective charges $\Zst_{\kappa}$, and the electronic
dielectric tensor $\epsinf$; \textsc{epw} supplies the e--ph coupling in the localized
Wannier representation $g(\Rv_{e},\Rv_{p})$, from which $g(\kv,\qv,\nu)$ is recovered by an
inverse Fourier transform.

\paragraph{Why LiF.}
Rocksalt LiF has a single, non-degenerate conduction-band minimum at $\Gamma$, one
dominant longitudinal-optical (LO) branch, and an intermediate-to-strong Fr\"ohlich
coupling $\alpha\approx4.8$. It is therefore the cleanest first material: nearly the
textbook Fr\"ohlich case, with a genuine analytic limit to validate against.

\paragraph{Units.}
The ab-initio model works in \emph{physical} units --- energies in eV, imaginary time in
eV$^{-1}$. The electron dispersion $\varepsilon(\kv)$ is the real Wannier-interpolated
conduction band from $H(\Rv)$ (not a parabola), and the phonons are the real branches
$\omega_{\qv\nu}$ from the polar-restored dynamical matrix, with \emph{all} modes retained.
The Fr\"ohlich coupling is still computed from the macroscopic constants as a cross-check,
\begin{equation}
  \alpha \;=\; \frac{e^{2}}{\hbar}\,
  \sqrt{\frac{m^{*}}{2\hbar\omega_{\mathrm{LO}}}}\,
  \left(\frac{1}{\varepsilon_{\infty}}-\frac{1}{\varepsilon_{0}}\right)
  \;\approx\; 4.76\quad(\text{LiF}),
  \label{eq:alpha}
\end{equation}
which fixes the weak-coupling target $E^{(2)}\!\approx\!-\alpha\,\hbar\omega_{\mathrm{LO}}$
against which the vertex scale is checked (Section~\ref{sec:pitfalls}, P8).

\section{Goals}
\label{sec:goals}

\begin{description}
  \item[G1 --- Real vertex.] Drive the DiagMC sampler with the \emph{ab-initio} e--ph
    vertex in place of the analytic one. The vertex is decomposed as
    \begin{equation}
      g_{mn\nu}(\kv,\qv) \;=\; \gshort_{mn\nu}(\kv,\qv)\;+\;\glong_{mn\nu}(\kv,\qv),
      \label{eq:split}
    \end{equation}
    the short-range (Wannier-interpolated) part plus the long-range polar part.

  \item[G2 --- Generic engine, native model.] Factor the sampler into a state-generic
    \texttt{rmc-diagmc} engine (\texttt{Runner}, \texttt{Update}, \texttt{MetropolisKernel},
    decay estimator) behind a \texttt{Model} trait. The ab-initio model is a \emph{native}
    implementation of that trait --- real $\varepsilon(\kv)$, $\omega_{\qv\nu}$, complex
    $g(\kv,\qv,\nu)$, and a $|g|^{2}$-weighted phonon-momentum proposal --- \emph{not} a
    reweighting of the Fr\"ohlich envelope. The Fr\"ohlich model is a second implementation of
    the same trait, which validates the engine at stage~0 (Section~\ref{sec:validation}).

  \item[G3 --- Correct polar long range.] Restore the long-range Fr\"ohlich
    $1/q$ divergence that dominates the LiF polaron. \textsc{epw} was run with
    \texttt{lpolar=.true.}, which \emph{subtracts} the long-range part before
    Wannierization --- from \emph{both} the e-ph vertex and the dynamical matrix --- so
    $\gshort$ from the tensor is short-ranged and both the vertex term $\glong$
    (Eq.~\eqref{eq:longrange}) and the dynamical-matrix term $\Dlong$
    (Eq.~\eqref{eq:dynlong}) must be added back analytically. The latter is what restores
    the physical TO frequency and the LO--TO splitting; without it the phonon zero-point
    factor $(2M_{\kappa}\omega)^{-1/2}$ in $\glong$ is mis-scaled (Section~\ref{sec:pitfalls},
    P7). \emph{Done:} both implemented as from-scratch ports of \textsc{epw}'s
    \texttt{rgd\_blk\_epw} and \texttt{rgd\_blk} (Ewald parameters $\Lambda_{\mathrm{Ew}}=1$,
    $g_{\max}=14$, $e^{2}=2$~Ry) and validated to machine precision against the coarse Bloch
    \texttt{.epb} and dynamical matrix \texttt{dynq} (Section~\ref{sec:pitfalls}, P6--P7).

  \item[G4 --- Clean Rust reader.] Implement the \textsc{epw}-file readers and the Wannier
    reconstruction from scratch in Rust, with no Python at runtime. The existing Python
    post-processor (\texttt{patt126-ab-initio-diagmc}) is used as a \emph{reference only},
    not copied; the rewrite also fixes a latent index assumption in the reference.

  \item[G5 --- Two-front validation.] Benchmark the LiF electron polaron against
    \emph{both}
    (i) Fr\"ohlich-limit DiagMC energetics at $\alpha_{\mathrm{LiF}}$
        (Mishchenko \emph{et al.}, 2000), and
    (ii) \emph{ab-initio} LiF polaron numbers
        (Sio, Verdi, Ponc\'e, Giustino, 2019).
\end{description}

\noindent
The long-range term used for G3 (Verdi--Giustino form) is
\begin{equation}
  \glong_{mn\nu}(\kv,\qv) \;=\;
  i\,\frac{4\pi e^{2}}{\Omega}\sum_{\kappa}
  \left(\frac{\hbar}{2M_{\kappa}\omega_{\nu\qv}}\right)^{\!1/2}
  \sum_{\Gv\neq-\qv}
  \frac{(\qv+\Gv)\cdot\Zst_{\kappa}\cdot\mathbf{e}_{\kappa\nu}(\qv)}
       {(\qv+\Gv)\cdot\epsinf\cdot(\qv+\Gv)}\,
  e^{-i(\qv+\Gv)\cdot\bm{\tau}_{\kappa}}\,
  e^{-\frac{(\qv+\Gv)\cdot\epsinf\cdot(\qv+\Gv)}{4\Lambda}}
  \bigl[U_{\kv+\qv}U_{\kv}^{\dagger}\bigr]_{mn},
  \label{eq:longrange}
\end{equation}
with $\Omega$ the cell volume, $M_{\kappa},\bm{\tau}_{\kappa}$ the atomic masses and
positions, $\mathbf{e}_{\kappa\nu}$ the phonon eigenvectors, and $\Lambda$ an Ewald
range-separation parameter. The final factor $[U_{\kv+\qv}U_{\kv}^{\dagger}]_{mn}$ is the
electron-band rotation into the smooth (Wannier) gauge; at the $\Gamma$-slice ($\kv=\Gamma$,
single conduction band) it reduces to the conduction-band Bloch overlap
$\langle u_{\Gamma+\qv}|u_{\Gamma}\rangle$. Restoring it is essential to the exact
short/long decomposition of Eq.~\eqref{eq:split}; see Section~\ref{sec:pitfalls}, P6. As
$q\to0$ the term behaves as $1/q$, reproducing the Fr\"ohlich tail
$|g_{F}(\qv)|^{2}\propto\alpha/q^{2}$.

\noindent
The \emph{same} dipole physics enters the dynamical matrix: \textsc{epw} likewise removes
the rigid-ion (dipole--dipole) long-range from $D$ before Wannierization, so the
interpolated phonons are short-ranged --- the optical triplet stays degenerate at $\Gamma$
at an unphysical frequency with no LO--TO splitting. The term to add back (Gonze form,
\textsc{epw} \texttt{rgd\_blk}) is, with $\Qv\equiv\qv+\Gv$,
\begin{equation}
  \Dlong_{\kappa\alpha,\kappa'\beta}(\qv) \;=\; \frac{4\pi e^{2}}{\Omega}
  \sum_{\Gv}
  \frac{ e^{-\Qv\cdot\epsinf\cdot\Qv/4\Lambda} }{ \Qv\cdot\epsinf\cdot\Qv }\,
  (\Qv\!\cdot\!\Zst_{\kappa})_{\alpha}\,(\Qv\!\cdot\!\Zst_{\kappa'})_{\beta}\,
  e^{i\Qv\cdot(\bm{\tau}_{\kappa}-\bm{\tau}_{\kappa'})}
  \;-\; \delta_{\kappa\kappa'}\,S_{\kappa\alpha\beta},
  \label{eq:dynlong}
\end{equation}
where the self-term $S_{\kappa\alpha\beta}=\sum_{\kappa''}\Dlong_{\kappa\alpha,\kappa''\beta}(\mathbf{0})$
(the same sum at $\qv=0$, summed over the second atom) is subtracted on the diagonal
blocks to enforce the acoustic sum rule. Diagonalizing $D^{\mathcal{S}}+\Dlong$
(mass-reduced) gives the physical branches $\omega_{\qv\nu}$ and eigenvectors
$\mathbf{e}_{\kappa\nu}(\qv)$ that enter Eq.~\eqref{eq:longrange}.

\section{Non-goals (deliberate scope of the first prototype)}
\label{sec:nongoals}

Several early non-goals of the \emph{first} prototype have since been \textbf{lifted}: the
$\Gamma$-slice vertex, the single Einstein branch, the ``no tensor compression'' rule, the
reweighted sampler, the Fr\"ohlich-unit engine constants, and the ``no GPU path'' rule
(N7). The model now uses the full electron-momentum vertex $g(\kv,\qv,\nu)$, all phonon
branches, an SVD-compressed Wannier tensor evaluated on a grid, a native $|g|^{2}$ momentum
proposal, physical units, and a validated CUDA backend (Section~\ref{sec:gpu}). What
remains deliberately out of scope:

\begin{description}
  \item[N1 --- Electron polaron only.] The single, non-degenerate conduction-band minimum
    at $\Gamma$. The hole/valence polaron (degenerate top valence band) is out of scope.

  \item[N8 --- Single material.] LiF is the sole validation target; this is not yet a
    general multi-material framework.
\end{description}

\noindent
(The retired non-goals keep their original labels N2--N7 where referenced elsewhere; they
are now realized features, recorded as done in Section~\ref{sec:deferred}.)

\section{Deferred to future work}
\label{sec:deferred}

\paragraph{Done.} The full-$\kv$ vertex with SVD compression and on-grid evaluation (former
N2/N4), all dispersive branches $\omega_{\nu}(\qv)$ (N3), the non-parabolic Wannier
dispersion $\varepsilon(\kv)$ via the generic \texttt{Model} trait (N6), the native
$|g|^{2}$ momentum sampler (N5), and the CUDA backend (former N7/D3,
Section~\ref{sec:gpu}) are now \emph{implemented}.

\paragraph{Remaining.}
\begin{description}
  \item[D1 --- Grid convergence.] \emph{Done} for the energy. The vertex is tabulated on a uniform $n^{3}$ grid; the discrete BZ sum underweights the near-$\Gamma$ $1/q^{2}$ polar region, so the polaron energy converges from above with a leading $1/n$ error (at $n{=}8$, $E_{0}\approx-0.122$~eV). Ladders $n=20\ldots120$ at $\tau=60$ and $\tau=240$~eV$^{-1}$ and their extrapolation are in Section~\ref{sec:discretization}. Still open: the $\tau{=}480$ guard at $120^{3}$, and a $\Gamma$-cell (Fr\"ohlich control-variate) correction that would remove the leading $1/n$ term rather than fit it.

  \item[D2 --- Hole polaron.] The degenerate valence-band-maximum polaron (lifting N1),
    needing a multi-band, degenerate (matrix-product sign-fix) treatment.

  \item[D4 --- More materials.] Generalize beyond LiF once the LiF prototype is validated
    (lifting N8).
\end{description}

\section{Validation stages}
\label{sec:validation}

The goals of Section~\ref{sec:goals} are checked in stages of increasing physical content:
\begin{enumerate}
  \item \textbf{Engine (stage 0).} The analytic Fr\"ohlich model, run through
    \texttt{rmc-diagmc} and the decay estimator, reproduces Mishchenko \emph{et al.} (2000):
    $E_{0}(\alpha{=}1)=-1.013$ and $E_{0}(\alpha{=}2)=-2.08$ within $3\sigma$, sign
    $\equiv1$. This validates the sampler, estimator, and sign tracking independently of the
    ab-initio vertex.
  \item \textbf{Vertex scale.} The ab-initio second-order energy $E^{(2)}(\Gamma)$ recovers
    the Fr\"ohlich weak-coupling target $-\alpha\,\hbar\omega_{\mathrm{LO}}$ (up to the
    coarse-grid underestimate), confirming the $(2M_{\kappa}\omega)^{-1/2}$ vertex
    normalization (Section~\ref{sec:pitfalls}, P8).
  \item \textbf{Compression convergence.} The all-orders $E_{0}$ is certified against the
    compression \emph{in the sampler itself}: the production rank budget must agree with the
    uncompressed vertex to within a stated tolerance,
    $\Delta_{\mathrm{rank}}=E_{0}(\text{budget})-E_{0}(\text{uncompressed})$, and the measured
    $\Delta_{\mathrm{rank}}$ is quoted as a systematic on every production number
    (Section~\ref{sec:discretization}).
  \item \textbf{Material comparison.} The grid-ladder formation energy against FEP-DMC (Luo, Park, Bernardi, 2025): their $E_{0}(20^{3})=-0.2479$~eV (from their example output, a single run at their effective $\tau\approx232$~eV$^{-1}$) and $-0.408$~eV in the $1/n\to0$ limit (linear fit through their three largest grids). At matched $\tau=240$~eV$^{-1}$ we measure $E_{0}(20^{3})=-0.2506\pm0.0003$~eV, $2.7$~meV deeper; the extrapolated $E_{\infty}=-0.4063\pm0.0006\,(\text{stat})\pm0.0008\,(\text{fit model})$~eV is $1.7$~meV shallower (Section~\ref{sec:discretization}). The two ladders therefore differ in slope by a few meV over $n=20\to\infty$, plausibly from dataset provenance (pseudopotentials, cutoffs, their rank budget); agreement is at the $1\%$ level, and a meV-level claim is not supported. Neither number includes the remote-band (Allen--Heine--Cardona) correction, which FEP-DMC adds separately to reach $-0.538$~eV.
\end{enumerate}

\section{Discretization, compression, and the thermodynamic limit}
\label{sec:discretization}

The object the sampler actually solves is not the continuum polaron Hamiltonian but its
restriction $\hat H_{n}$ to a uniform $n^{3}$ Brillouin-zone grid shared by electrons and
phonons. The grid is closed under momentum addition, so every diagram's momentum
bookkeeping stays on it; the per-arc momentum measure is the lattice normalization
$\frac{1}{N_{q}}\sum_{\qv\neq\Gamma}$ with $N_{q}=n^{3}$ (matching the standard
first-principles self-energy convention), and the DiagMC solves each $\hat H_{n}$
\emph{numerically exactly}. All physical statements are then statements about the sequence
$E_{0}(n)$: because the discrete sum underweights the $1/q^{2}$ polar region within a
$1/n$-neighbourhood of $\Gamma$, $E_{0}(n)$ converges from above with a leading error
linear in $1/n$, and the thermodynamic limit is read off a $1/n\to0$ extrapolation
over a ladder of grids --- the same methodology as FEP-DMC, whose
LiF benchmarks are $E_{0}(20^{3})=-0.248$~eV and $-0.408$~eV extrapolated.

\paragraph{The LiF ladder and its error budget.} Banked rungs: $\tau=60$~eV$^{-1}$ at $n=20,32,40,50,60,80,100,120$ and $\tau=240$~eV$^{-1}$ at $n=20,40,60,80,100,120$ (runbook \texttt{notebooks/LiF-e0-gamma-extrapolation.ipynb}). Only the leading $u=1/n$ term has a physical origin (the missing $\Gamma$ cell), so the model $E_{0}(u)=E_{\infty}+Cu+Du^{2}$ is fitted by weighted least squares, and the spread over fit models and windows is quoted as a separate systematic. For $\tau=240$ the quadratic fit over all rungs gives $E_{\infty}=-0.4063\pm0.0006$~eV ($\chi^{2}_{\mathrm{red}}=1.4$). Acceptable alternatives ($\chi^{2}_{\mathrm{red}}\le2.2$, own error $\le2$~meV: linear from $n\ge40$, linear through the three largest grids as in FEP-DMC, quadratic from $n\ge40$) span $-0.4064\ldots-0.4048$~eV, i.e.\ a fit-model systematic of $\pm0.8$~meV. Underdetermined fits (cubic on five points, own error $\approx6$~meV) are uninformative rather than competing models, and are excluded. An independent route agrees: the $\tau$ shift $E_{0}(\tau240)-E_{0}(\tau60)$ is $n$-independent for $n\ge40$ ($-2.83\pm0.13$~meV), and adding it to the better-conditioned eight-rung $\tau60$ quadratic fit ($-0.4030\pm0.0002$~eV) gives $-0.4058$~eV. The compression systematic $\Delta_{\mathrm{rank}}(800)=+0.5\pm0.3$~meV (at $20^{3}$) is below the fit-model systematic. The $\tau$ systematic is certified by a $\tau=480$~eV$^{-1}$ guard at $n=120$ (pending). Current quote: $E_{\infty}=-0.4063\pm0.0010$~eV (stat.\ and fit model combined), before the $\tau$ guard.

\paragraph{Compression as a controlled, quoted approximation.}
The short-range vertex is compressed to a global rank budget (the largest singular values
across all band pairs); the long-range polar term is always exact. The acceptance
criterion is \emph{never} element-wise reconstruction error --- it is convergence of the
physical observable, certified in the sampler itself:
\[
  \Delta_{\mathrm{rank}} \;=\; E_{0}(\text{budget}) - E_{0}(\text{uncompressed}),
\]
measured by two otherwise-identical DMC runs and quoted as a systematic error alongside
the statistical one on every production result. Two empirical lessons fixed this
protocol. First, the cheap proxy --- the second-order energy $E^{(2)}$ at $\kv=\Gamma$ ---
systematically \emph{overestimates} the true sampler systematic by a factor of $3$--$10$,
because the sampler averages the compression error over the whole zone while the proxy
probes a single $\kv$-row; the proxy is useful for scanning, but only the in-sampler
measurement decides. Second, convergence with the global budget is \emph{non-monotonic}
(singular values rank global tensor weight, not importance for a given observable), so
``two nearby budgets agree'' is not evidence of convergence --- only the comparison
against the uncompressed vertex is. For the current LiF dataset the certified production
budget is $800$ of $1600$ components with $\Delta_{\mathrm{rank}}=1.0\pm0.8$~meV; the
FEP-DMC-sized budget ($\approx300$) measured a genuine $4.6\pm0.8$~meV and was rejected.

\paragraph{Dataset provenance.}
The DFPT/Wannier inputs use $8\times8\times8$ coarse $k$- and $q$-grids, matching FEP-DMC's
LiF setup. Upgrading from the initial $6\times6\times6$ deepened $E_{0}(8^{3})$ by
$\sim11$~meV --- a pure interpolation-quality correction, worth recording because it is the
same order as the compression and statistical errors being controlled above. One physical
subtlety appeared with the finer grid: it contains the $W$ star, where the two lowest
conduction bands are exactly degenerate, making per-band comparisons gauge-ambiguous
there (the round-trip validation skips those points; the sampler itself is unaffected,
since it consumes gauge-invariant $|g|^{2}$ combinations).

\section{Performance: the ab-initio vertex cache}
\label{sec:performance}

The full electron-momentum vertex --- the former $\Gamma$-slice non-goal (N2), now lifted ---
cannot be reduced to the scalar $|g(\Gamma,\qv,\mathrm{LO})|^{2}$. Instead the Wannier e--ph tensor is
compressed by a per-band-pair truncated SVD into rank-$r$ factors $u_{g}(\kv)$,
$v_{g}(\qv,\nu)$, tabulated once on the $q$-grid; a vertex evaluation is then the rank-$r$
contraction
\[
  g(\kv,\qv,\nu)\;=\;\sum_{g=1}^{r} s_{g}\,u_{g}(\kv)\,v_{g}(\qv,\nu)\;+\;\glong(\kv,\qv,\nu),
\]
called millions of times per DiagMC chain, so its cost sets the wall-clock. On an $8^{3}$
grid at $r\!\sim\!150$ the factor tables are $\approx0.18$~GiB --- far larger than cache ---
so the memory layout of the component axis $g$, not the arithmetic, dominates.

\paragraph{Cache-hostile layout (before).} Storing the factors component-major,
$u[g][\kv]$ and $v[g][\qv][\nu]$, places successive components of a single vertex sum
$\sim\!49$~KB apart, so each of the $r$ terms is a separate cache miss and the inner
contraction is bound entirely by memory latency.

\paragraph{Fix: transpose the component axis inward.} Storing the factors $\kv$-major /
$(\qv,\nu)$-major --- $u[\kv][g]$, $v[\qv][\nu][g]$ --- makes the $r$ components of one
vertex evaluation contiguous, so the contraction streams two cache-line runs instead of
striding $r$ scattered lines. The precompute is reorganized to parallelize over $\kv$
(rather than over $g$) so it writes this layout directly.

\paragraph{Effect.} For the $n_{\mathrm{sv}}$ convergence sweep ($500$k steps $\times$ $10$
chains, $8^{3}$ grid), a single rank drops from $258\to23$~s at $r{=}50$ and
$1033\to54$~s at $r{=}100$ --- an $11$--$19\times$ wall-clock reduction that grows with $r$.
The evaluation is now memory-bandwidth bound on the factor tables; the $O(\text{order})$
phase recomputation in the decay estimator is a smaller secondary term, left for later.

\section{GPU sampling of the discretized Hamiltonian}
\label{sec:gpu}

The former non-goal N7 is lifted: the ab-initio engine runs on CUDA (CubeCL), and the port
exists to \emph{win on throughput}, unlike the earlier Fr\"ohlich GPU port (a correctness
vehicle). The physics favors the GPU here because the dominant cost is not pointer chasing
but the vertex evaluation of Section~\ref{sec:performance}: a contiguous rank-$r$ complex
dot streaming $\approx16r$ bytes per call --- a memory-\emph{bandwidth} problem, and HBM
bandwidth is the one resource where a GPU outclasses a CPU node by an order of magnitude.
Measured on an A100 at production rank ($r=800$, $50^{3}$ grid), the cooperative evaluation
sustains $\approx1.24$~TB/s ($\approx80\%$ of HBM peak; $\approx90\%$ at $20^{3}$ where L2
helps), versus $14$--$76$~GB/s for the same evaluation on CPU nodes.

\paragraph{Parallelism and the cooperative vertex dot.}
Chains never communicate, so the natural mapping is many independent chains per device.
The microbenchmark that gated the port showed that \emph{per-thread} evaluations never
coalesce ($0.4\times$ a CPU node), so the kernel assigns \textbf{one chain per 32-lane
warp}: every lane executes the chain's update logic \emph{redundantly} --- the RNG is
counter-based (Philox4x32-10 keyed by $(\text{seed},\text{chain},\text{step})$), so all
lanes compute identical draws with no communication, control flow is warp-uniform by
construction, and lanes cooperate only inside vertex evaluations, striding the rank axis
with a shared-memory reduce. One arc factor $g(\kv_{e},\qv,\nu)\,g^{*}(\kv_{a},\qv,\nu)$
fuses both dots over a single shared $v_{g}(\qv,\nu)$ strip, halving the dominant table
traffic. Diagram state is a dense $\tau$-sorted array per chain (positional keys mirroring
the CPU engine one-to-one); one lane applies accepted mutations between barriers. Tables
are stored planar (separate re/im f32), $\tau$ and all weight arithmetic are f64, and
momenta are exact integer grid indices, as on the CPU.

\paragraph{Why lockstep execution is mandatory, and why it does not bias the sampling.}
The warp-uniformity above is a hardware \emph{correctness} requirement, not a tuning
choice. The cooperative reduce synchronizes the warp with block-wide barriers, and the
hardware pairs barrier arrivals by \emph{dynamic count}, not by source location: had lanes
diverged, either some would wait at a barrier the others never reach (deadlock), or the
counts would happen to match at \emph{different} program points and lanes in unrelated
phases would race through the one shared scratch buffer (silent bit-level corruption).
Since the number of barriers a step executes depends on the update type, the proposal
draws, and the diagram state, control flow must agree across lanes over the \emph{entire
dynamic trace} --- which redundant execution of identical draws guarantees bit-for-bit.
The statistical side is equally clean. The update-\emph{type} schedule is drawn from a
dedicated counter-based stream keyed by (seed, chain, step) --- never by the state.
Because the schedule is independent of the state and of the chain's proposal/acceptance
stream, each chain is exactly the standard random-scan chain driven by the mixture kernel
$n_{\mathrm{upd}}^{-1}\sum_{u}P_{u}$, and the $1/n_{\mathrm{upd}}$ selection probabilities
of paired moves (phonon insertion $\leftrightarrow$ removal) cancel in the acceptance
ratio. The load-bearing invariant is that selection stays state-independent: an
``impossible'' proposal (e.g.\ a removal at order zero) counts as proposed-and-rejected and is never
redrawn, since resampling until a possible move would make the effective selection
distribution state-dependent and break the cancellation --- a genuine bias, not a
determinism nit.

\paragraph{Determinism and multi-GPU scaling.}
Because chain streams are keyed by \emph{global} chain ids, a population split across
ranks (one rank per GPU, chain offset $=\text{rank}\times C$) is a pure repacking of a
single larger run: per-chain measurement batches concatenate in global chain order, and
the merged report is \textbf{bit-identical} to the equivalent single-device run ---
verified on a $2\times$A100 node ($2\times8192$ chains vs.\ one $16384$-chain run, every
estimator equal to the last bit). Multi-GPU throughput is therefore linear and the merged
statistics are a pure function of (seed, total chains), independent of packing.

\paragraph{Validation stages.}
Correctness is certified in stages, each mechanical to check: (i) a flat, fixed-capacity
CPU re-implementation of the four updates (the ``executable spec'') agrees with the
production engine's acceptance ratios to $10^{-10}$ on imported random diagrams; (ii) the
CUDA kernel, run at one lane per chain on CubeCL's CPU runtime, reproduces the flat driver
\emph{bit-for-bit} on the real LiF table (identical reports and accept counts); (iii) the
same kernel passes exact accept-count parity on real hardware at 32 lanes; (iv) the
physics anchor: the full $20^{3}$ ladder configuration on one A100 gives
$E_{0}=-0.24826\pm0.00023$~eV against FEP-DMC's $-0.2479$~eV. All floating-point
association in the kernel mirrors the CPU code term by term, so stage (ii) is exact rather
than statistical.

\paragraph{Throughput at matched physics.}
At $20^{3}$, rank $800$, $\langle n\rangle\approx16$, fixed $\tau=60$~eV$^{-1}$: one A100
sustains $8.79\times10^{6}$ steps/s (kernel-only; a full production-statistics $20^{3}$ ladder run
is $151$~s wall), versus $2.83\times10^{6}$ steps/s measured for the production CPU engine
on a 128-core Zen3 node with the same configuration --- $3.1\times$ per GPU, $6.2\times$
per two-GPU node. This is the CPU's \emph{best} case: the $20^{3}$ table ($0.34$~GB) fits
the node's aggregate L3, while at $50^{3}$ ($5.4$~GB) the CPU falls to DRAM bandwidth and
the GPU keeps streaming from HBM, so the margin grows with the grid --- exactly where the
D1 ladder needs it.

\paragraph{Estimator bookkeeping at GPU populations.}
Two scaling lessons from operating $10^{4}$ chains. First, the leave-one-out jackknife
must be computed as $\text{total}-\text{batch}$ (O(bins) per replicate); recomputing the
aggregate per replicate is $O(n_{\mathrm{batches}}^{2})$ and silently turns a two-minute
run into hours at $10^{5}$ batches. With chains independent, the natural batching is
\emph{one batch per chain}, jackknifing over chains. Second, warmup must exceed the
diagram-order equilibration transient ($\sim\!10^{5}$ steps for LiF at $20^{3}$)
\emph{regardless} of how many chains amortize the statistics: a many-short-chain ensemble
does not dilute the transient the way one long chain does. An early GPU run with $2\times
10^{4}$ warmup steps per chain came out $1.7$~meV shallow with $\langle n\rangle$ low by
$2\%$ --- statistically resolved at $5\sigma$ against the long-chain CPU value --- purely
from under-equilibration; production GPU configs now warm up $3\times10^{5}$ steps per
chain.

\section{Practical pitfalls and fixes}
\label{sec:pitfalls}

Turning the \textsc{epw} run into a \emph{usable} vertex, and then a trustworthy DiagMC
number, meant working through several failures whose symptoms were misleading. P1--P7 are
data-preparation (A0/A1); P8--P9 are the DiagMC consumer (Phase B). They are recorded here so
they are not rediscovered; the blow-by-blow A0/A1 diagnosis, with reproducing commands, lives
in \texttt{datasets/LiF/EPH-NAN-DIAGNOSIS.md}.

\begin{description}
  \item[P1 --- Ultrasoft pseudopotentials give an all-\texttt{NaN} e--ph tensor (root
    cause).]
    The first A0 deck used ultrasoft pseudopotentials (SSSP-efficiency GBRV). \textsc{epw}'s
    electron--phonon matrix elements are only supported with \emph{norm-conserving}
    pseudopotentials --- the ultrasoft augmentation charge is not carried into the coupling.
    The coupling came out finite \emph{only} at $\qv=\Gamma$ (the trivial reference $q$,
    which needs no symmetry rotation) and \texttt{NaN} at every other $\qv$; the
    Bloch$\to$Wannier sum then spread the \texttt{NaN} to $100\%$ of \texttt{LiF.epmatwp}.
    Because \texttt{ph.x} \emph{does} support ultrasoft, the phonons, dynamical matrix,
    $\Zst_{\kappa}$ and $\epsinf$ were all correct --- only the e--ph coupling was wrong,
    which is why the fault was easy to misattribute.
    \emph{Fix:} switch Li and F to norm-conserving Pseudo Dojo ONCV (scalar-relativistic
    PBE). This raises the required plane-wave cutoff
    ($\texttt{ecutwfc}:50\!\to\!96$~Ry, $\texttt{ecutrho}:400\!\to\!384$~Ry, i.e.\ the
    norm-conserving $4\times$), which makes \texttt{ph.x} $\sim\!3\text{--}4\times$ slower
    (the wavefunctions move from the small smooth FFT grid onto the dense grid) --- an
    unavoidable cost of the correct pseudopotential family. With norm-conserving pseudos the
    tensor is finite at every $\qv$, and $\epsinf$, $\Zst_{\kappa}$ are essentially unchanged.

  \item[P2 --- ``Right size, wrong values'': a silent failure.]
    The broken run finished cleanly and wrote \texttt{LiF.epmatwp} at exactly the expected
    byte size; the only hint was an easily-dismissed \texttt{IEEE\_INVALID} note at exit. The
    file was $100\%$ uninitialised-memory \texttt{NaN} (a uniform garbage word that even
    \emph{changed} between reruns). An earlier ``A0 done'' had been declared on file size
    alone, without ever checking the tensor \emph{values}.
    \emph{Fix:} the pipeline (\texttt{run\_pipeline.nu}) now runs an explicit finiteness
    check on \texttt{LiF.epmatwp} and aborts loudly on \texttt{NaN}, so a broken run can no
    longer masquerade as a good one.

  \item[P3 --- Two dead ends from assuming an upstream bug.]
    Before the pseudopotential cause was found, two ``\textsc{epw} bug'' hypotheses were
    pursued and \emph{both failed}: (i)~widening the Wannier outer window
    \texttt{dis\_win\_max} to enclose every band made the coarse tensor \emph{worse}, not
    better; (ii)~a bounds-guard patch to \textsc{epw}'s \texttt{gmap\_sym} routine (a genuine
    latent out-of-bounds read) compiled in cleanly yet left the output byte-identical --- so
    that read was not the cause. The lesson, recorded explicitly: suspect the \emph{inputs}
    (here, an unsupported pseudopotential family) before patching upstream code.

  \item[P4 --- The reference \texttt{.epb} reader has the wrong record layout.]
    The Python reference (\texttt{patt126}) reads the coarse Bloch \texttt{.epb} file
    assuming a single band count throughout the record. The real record interleaves a
    \emph{20-band} eigenvalue array with a \emph{15-band} coupling array (of which only a
    $5\times5$ block is populated); a na\"ive parse silently mis-aligns and yields a
    \emph{false positive} --- a ``the coarse data looks fine'' reading that briefly sent the
    debugging down a wrong path.
    \emph{Fix:} the true record layout was reverse-engineered and byte-verified --- a
    \emph{single} unformatted Fortran record (4-byte length markers)
    $[\,n_{qc},\,\mathbf{x}_{qc}(3,n_{qc}),\,\varepsilon(20,n_{k}),\,D(6,6,n_{qc}),\,
    g(15,15,n_{k},6,n_{qc}),\,\Zst,\epsinf\,]$, all little-endian and column-major, one file
    per pool with the pool-local $n_{k}$ recovered from the record length. Two indexing
    traps: the coupling's mode axis sits \emph{interior} (stride $n_{k}\!\cdot\!6$, not
    $6\!\cdot\!n_{k}$), and the conduction-band minimum is kept-band local index $0$ in the
    coupling array but eigenvalue index $5$ in $\varepsilon$. The Rust reader uses this
    directly and decodes only the needed conduction-band slice on demand.

  \item[P5 --- \textsc{epw} coarse-grid prerequisites (reproducibility).]
    Several setup constraints are load-bearing and non-obvious: the \texttt{nscf}
    $k$-points must be an \emph{explicit} $\Gamma$-centred crystal list (an automatic,
    half-shifted Monkhorst--Pack mesh breaks \textsc{epw}'s $\kv+\qv$ folding); the whole
    \texttt{scf/ph/nscf/epw} chain must use \emph{one} Quantum \textsc{espresso} version (a
    mixed-version \texttt{phsave} corrupts the heap); and \textsc{epw} needs a large stack
    (\texttt{ulimit -s unlimited}). Finally, changing the pseudopotentials shifts the band
    edges ($\sim\!0.15$~eV here), so the disentanglement windows (\texttt{dis\_froz\_max},
    \texttt{dis\_win\_max}) must be re-tuned against the new \texttt{nscf} eigenvalues ---
    otherwise the frozen conduction-band minimum, which must stay exact, can slip out of its
    window.

  \item[P6 --- The polar term lives in the Cartesian displacement basis, and carries an
    electron-band rotation.]
    Two subtleties in Eq.~\eqref{eq:longrange} decide whether the split
    $g=\gshort+\glong$ of Eq.~\eqref{eq:split} actually closes, and both were found by
    round-tripping against the coarse \texttt{.epb} (which stores the \emph{full},
    pre-subtraction coupling).
    \emph{(i)~Basis.} \textsc{epw} subtracts $\glong$ from the coupling in the
    \emph{Cartesian atomic-displacement} basis ($\kappa\alpha$, i.e.\ per atom and Cartesian
    direction), \emph{before} the phonon-eigenvector projection $\mathbf{e}_{\kappa\nu}(\qv)$
    and the $(\hbar/2M_{\kappa}\omega)^{1/2}$ mode factor. So the reconstructed short-range
    part and the re-added $\glong$ must both be formed in that basis and summed there;
    projecting onto branches first is inconsistent. The mode factor and eigenvector
    projection are applied afterwards by the consumer.
    \emph{(ii)~Gauge.} The band-rotation factor $[U_{\kv+\qv}U_{\kv}^{\dagger}]_{mn}$ --- at
    the $\Gamma$-slice the conduction-band Bloch overlap
    $\langle u_{\Gamma+\qv}|u_{\Gamma}\rangle$ --- is easy to drop under the ``vertex depends
    on $\qv$ only'' reduction (both electron legs rotated with $U_{\Gamma}$, overlap
    $\equiv1$). Dropping it puts the polar term in the wrong electron gauge: near $\Gamma$
    the reduced sum still overshoots $\lVert g\rVert$ by $\sim\!9\%$, and the overlap itself
    falls from $\approx0.94$ at the first coarse shell to $0$ at the zone boundary as the
    conduction-band character changes with $\kv$. Restoring the exact bra ($u_{\Gamma+\qv}$)
    and the overlap makes the decomposition exact.
    \emph{Result:} against all $215$ non-zero coarse $\qv$,
    $g_{\mathrm{full}}=\gshort+\langle u_{\Gamma+\qv}|u_{\Gamma}\rangle\,\glong$ to machine
    precision --- gauge-invariant cosine $=1.000000$ and $|\text{ratio}-1|\sim3\times10^{-7}$
    (the \texttt{.epb} Kohn--Sham states carry an arbitrary per-$\kv$ phase, so the check is
    on direction and magnitude, not the raw complex difference). The reduced $\qv$-only
    vertex remains the object the sampler consumes; the exact bra $+$ overlap is used only to
    prove the reconstruction and the polar term against real data.

  \item[P7 --- The dynamical matrix carries the same polar long-range.]
    The subtraction of P6 is not confined to the e-ph vertex: \textsc{epw} removes the
    rigid-ion dipole--dipole term from the dynamical matrix as well (\texttt{rgd\_blk}), so
    the Wannier-interpolated phonons are short-ranged --- the optical triplet stays
    degenerate at $\Gamma$ at an unphysical $11.95$~THz with no LO--TO splitting, where the
    physical TO is $8.12$~THz (agreeing with \texttt{ph.x}). This is easy to overlook because
    the phonon \emph{eigenvectors} are barely affected (the optical modes are symmetry-fixed
    Li--F relative motion), but the \emph{frequency} enters the vertex through the zero-point
    factor $(2M_{\kappa}\omega_{\qv\nu})^{-1/2}$, so a short-range $\omega$ mis-scales
    $|g|^{2}$ by a factor of order two.
    \emph{Fix:} port \textsc{epw}'s \texttt{rgd\_blk} (Gonze form, Eq.~\eqref{eq:dynlong}) --
    the $\qv\neq0$ dipole--dipole sum over all atom-pair blocks plus the $\qv{=}0$ self-term
    on the diagonal blocks -- and add it back before diagonalizing.
    \emph{Validated} at two levels: the reconstructed long-range term equals coarse-%
    \texttt{dynq} minus the short-range force constants \emph{element-by-element} (ratio
    $1.000$), and all six branches then reproduce the DFPT \texttt{dynq} at the coarse $q$
    ($[2.76,2.76,5.23,8.12,8.12,18.45]$~THz, LO $=18.45$) --- the phonon analog of the P6
    round-trip. The restored $\omega_{\mathrm{LO}}(\Gamma)=18.67$~THz gives a
    Lyddane--Sachs--Teller ratio $\varepsilon_{0}/\varepsilon_{\infty}=
    (\omega_{\mathrm{LO}}/\omega_{\mathrm{TO}})^{2}\approx5.3$, i.e.\ $\varepsilon_{0}\approx11$
    (experiment $\approx9$).

  \item[P8 --- The phonon quantization factor was dropped in the branch projection
    (Phase B).]
    Distinct from P7 (where the \emph{frequency} was wrong): here $\omega_{\qv\nu}$ was
    correct, but the consumer projected the Cartesian coupling onto phonon branches with only
    the eigenvector $\mathbf{e}_{\kappa\nu}$ and the Ry$\to$eV conversion, \emph{omitting} the
    $(2M_{\kappa}\omega_{\qv\nu})^{-1/2}$ zero-point factor of Eq.~\eqref{eq:longrange}
    entirely. Every $|g|$ was then $\sim\!7\times$ too large, so the second-order energy came
    out $\approx-11$~eV ($\sim\!30\times$ too deep) and the DiagMC diagram order ran away (no
    order-$0$ samples, the estimator's normalization sector). The bug hid because the
    \emph{oracle} used for the SVD cross-check made the same omission, and the \texttt{.epb}
    test validates only the \emph{Cartesian} coupling, above the projection.
    \emph{Fix:} apply $(2M_{\kappa}\omega_{\qv\nu})^{-1/2}$ in the branch projection of
    \emph{both} the tabulated vertex and the direct oracle, guarding the $\omega\to0$ acoustic
    modes at $\Gamma$. \emph{Validated:} $E^{(2)}$ fell to $-0.10$~eV on the $8^{3}$ grid ---
    the expected $\sim\!3.5\times$ coarse-grid underestimate of the continuum
    $-\alpha\hbar\omega_{\mathrm{LO}}\approx-0.37$~eV --- and the sampler then equilibrates at
    order $\approx3$.

  \item[P9 --- Reading $E_{0}$ from the decay estimator: $\mu$ and $\tau_{\max}$ (Phase B).]
    The ground state is the large-$\tau$ slope of $G(\tau)\sim Z\,e^{-(E_{0}-\mu)\tau}$, and
    two knobs must bracket it. The chemical potential $\mu$ must sit \emph{just below}
    $E_{0}$: if $\mu>E_{0}$ the weight grows with $\tau$ and the diagram order runs away
    (symptom: ``no order-$0$ samples''); if $\mu$ is far below $E_{0}$ the phonon cloud is
    starved and the fit window falls where $G$ is exponentially small, so $E_{0}$ scatters
    with underestimated error bars. The window must also satisfy
    $\tau_{\mathrm{fit}}\gg1/(E_{1}{-}E_{0})\sim1/\omega_{\mathrm{LO}}$ to clear
    excited-state contamination. The acceptance criterion is that the extracted
    $E_{0}=\mu+\text{slope}$ is \emph{$\mu$-independent}. On the coarse $8^{3}$ grid the
    polaron is effectively strong-coupling (order $3$--$5$, $E_{0}\approx-0.26$~eV), so
    $\mu\approx-0.5$~eV with $\tau_{\max}\approx30$~eV$^{-1}$ is the working regime.
\end{description}

\section{The finite-$\beta$ thermodynamic energy estimator (FEP-DMC)}
\label{sec:thermo-estimator}

The energy in Section~\ref{sec:pitfalls} (P9) is read from the large-$\tau$ decay
of $G(\tau)$, which carries a documented sensitivity to the fit window $\tau_{\max}$
at strong coupling. FEP-DMC (Luo, Park, Bernardi, \emph{Nat.\ Phys.}\ \textbf{21}, 1275
(2025); code a Perturbo fork) instead measures the polaron energy as a
\emph{thermodynamic} (operator) expectation value on each worldline, with \emph{no}
fit. This section records how that estimator works and its single-band form. It is
\textbf{implemented} as the production estimator: the sampler runs in a fixed-$\tau$
ensemble (external time pinned at $\tau_{\max}$), where the chemical potential $\mu$
drops out of the ensemble entirely --- eliminating the $\mu$-bracketing of P9 --- and the
local energy $E_{\mathrm{loc}}=\tau^{-1}\bigl[\sum_{\mathrm{seg}}\varepsilon\,\Delta\tau
+\sum_{\mathrm{arcs}}\omega\,\Delta\tau-2N_{\mathrm{arcs}}\bigr]$ is averaged with the
sign. It is validated against an exactly solvable dispersionless (Lang--Firsov) lattice
model, $E_{0}=-\bar{g}^{2}/\omega$ at all couplings, and against the Fr\"ohlich
$\alpha=1,2$ benchmarks, where it agrees with the decay fit at lower variance.

\paragraph{Free-$\tau$ vs.\ fixed-$\tau$ ensembles.}
Both estimators sample diagrams of the zero-temperature imaginary-time Green's function
\begin{equation}
  G(\mathbf{p},\tau) \;=\;
  \langle 0|\,c_{\mathbf{p}}\,e^{-\tau(\hat H-\mu)}\,c_{\mathbf{p}}^{\dagger}\,|0\rangle
  \;\xrightarrow{\ \tau\to\infty\ }\; Z\,e^{-(E_{0}-\mu)\tau},
  \label{eq:green}
\end{equation}
a diagram being one electron backbone of external length $\tau$ dressed by phonon arcs
with internal vertex times and momenta. The two modes differ in whether that
\emph{external} $\tau$ is itself a sampled variable. In the \emph{free-$\tau$} ensemble
(\texttt{fixed\_tau: false}) the Markov chain updates $\tau\in(0,\tau_{\max})$ alongside
the internal degrees of freedom; histogramming the signed weights over $\tau$
(normalized through the analytically integrable order-zero sector) yields the whole
curve $G(\tau)$ in one run, from which the decay fit extracts both $E_{0}$ and the
quasiparticle weight $Z$, with the per-bin local-energy curve $E(\tau)$ as a built-in
convergence check. Its costs are those of Section~\ref{sec:pitfalls}, P9: statistics
diluted over the $\tau$ bins with only the large-$\tau$ tail informative, a genuine
fit-window systematic, and a load-bearing $\mu$ --- the weight carries $e^{\mu\tau}$,
which shapes the sampled $\tau$ distribution. In the \emph{fixed-$\tau$} ensemble
(\texttt{fixed\_tau: true}, FEP-DMC's \texttt{zeroTMC}; all production runs) the
external time is pinned at $\tau_{\max}$ and only internal degrees of freedom are
sampled: $e^{\mu\tau}$ is a constant that cancels in every ratio, so $\mu$ drops out
entirely, and all statistics sit at the single $\tau$ that matters. The energy then
follows from the internal-time rescaling identity: a diagram of $N$ arcs has $2N$
internal vertex times, so rescaling them by $\tau$ exposes the explicit
$\tau$-dependence $\tau^{2N}e^{-\sum\varepsilon\,\Delta\tau-\sum\omega\,\Delta\tau}$ of
each configuration's weight, whence
\begin{equation}
  -\partial_{\tau}\ln G(\tau)
  \;=\; \frac{1}{\tau}\Bigl\langle
    \sum_{\mathrm{seg}}\varepsilon_{i}\,\Delta\tau_{i}
    \;+\;\sum_{\mathrm{arcs}}\omega_{a}\,\Delta\tau_{a}
    \;-\;2N
  \Bigr\rangle_{\tau}
  \;=\; E_{0}-\mu+\mathcal{O}\bigl(e^{-\Delta\tau}\bigr),
  \label{eq:logderiv}
\end{equation}
the sign-weighted average taken in the fixed-$\tau$ ensemble. This is the
$E_{\mathrm{loc}}$ quoted above (the code's \texttt{local\_energy}), and equals
Eq.~\eqref{eq:estimator} in the $T\to0$ single-band limit up to the vertex-counting
convention ($-2N$ arcs here vs.\ $-(N_{v}-1)$ vertices there). The estimator therefore
approaches $E_{0}$ \emph{from above} at finite $\tau$, as $\sum_{i}w_{i}\Delta_{i}e^{-\Delta_{i}\tau}$ over excited states $i$ of the $\mathbf{p}$ sector with gaps $\Delta_{i}$ and relative weights $w_{i}$.

The relevant excited states at $\mathbf{p}=\Gamma$ are the polaron at $-\qv$ plus one phonon at $\qv$, with gap $\Delta(\qv)=E_{0}(-\qv)-E_{0}+\omega_{\nu}(\qv)$. For the LO branch, $\Delta\ge\omega_{\mathrm{LO}}=77.2$~meV (18.67~THz), which gives $\omega_{\mathrm{LO}}\tau\approx4.6$ at $\tau=60$~eV$^{-1}$ and $\approx18.5$ at $\tau=240$~eV$^{-1}$ (FEP-DMC's effective $1/k_{B}T$ at 50~K is $232$~eV$^{-1}$). The acoustic branches are gapless, however: on the $n^{3}$ grid the smallest gap is $\approx c_{s}q_{\min}\propto1/n$ ($4.4$~meV at $16^{3}$, below $1$~meV beyond $n\approx70$, where the soft-mode cutoff removes those modes). Convergence in $\tau$ is therefore \emph{not} uniform in $n$. The weak acoustic coupling keeps the amplitude small, but no single $\Delta\approx\omega_{\mathrm{LO}}$ controls it. Measured: $E_{0}(\tau240)-E_{0}(\tau60)=-0.9\pm0.3$~meV at $20^{3}$ but $-2.8\pm0.1$~meV (flat) for $n\ge40$. Because fixed-$\tau$ produces no curve whose flattening can be inspected, two further checks certify $\tau$ convergence: a $\tau=480$~eV$^{-1}$ guard at $n=120$, and a free-$\tau$ run at $40^{3}$. The free-$\tau$ run bins the same $E_{\mathrm{loc}}$ over $\tau\in(0,240]$; its per-bin average must reproduce the fixed-$\tau$ points at $\tau=60,240$, and its shape measures the effective gap directly.

\paragraph{Quantity measured.}
Rather than the asymptotic slope, the polaron energy is the operator average
\begin{equation}
  E_{\mathrm{pol}} \;=\; \langle \hat H\rangle
  \;=\; \frac{\operatorname{Tr}\!\bigl[\hat H\,e^{-\beta\hat H}\bigr]}
             {\operatorname{Tr}\!\bigl[e^{-\beta\hat H}\bigr]},
  \label{eq:thermo}
\end{equation}
evaluated on each Monte-Carlo worldline of total imaginary-time length $\tau=\beta$, where $\beta=1/k_{B}T$ is the code's \texttt{tauMax}. Here $\operatorname{Tr}$ is the matrix trace $\operatorname{Tr}[M]=\sum_{i}M_{ii}$; the trace over the complete basis is what turns Eq.~\eqref{eq:thermo} into a proper thermal (and, as $\beta\to\infty$, ground-state) expectation value. This is the standard Prokof'ev--Svistunov / Mishchenko thermodynamic estimator.

\paragraph{Worldline weight.}
A configuration is one electron backbone on $\tau\in[0,\beta]$, cut into segments by
vertices at $\tau_{1}<\tau_{2}<\cdots$, with phonon arcs joining vertex pairs. In
single-band form its weight is a product of
\begin{itemize}
  \item electron segments $s$ of duration $\Delta\tau_{s}$: $e^{-\varepsilon_{k_{s}}\Delta\tau_{s}}$
    (the code's \texttt{Gel}$(\varepsilon,\tau)=e^{-\varepsilon\tau}$ and \texttt{propagator}).
    The electron is kept at $T=0$; only phonons are thermal, as appropriate for a single carrier;
  \item phonon arcs $a$ of span $\Delta\tau_{a}$ at frequency $\omega_{a}$: the finite-$T$
    propagator $D(\omega,\tau)$ below;
  \item a vertex factor $g_{mn\nu}(\kv,\qv)$ at each of the $N_{v}$ vertices.
\end{itemize}
Thus $\ln w$ depends on $\beta$ through every segment length, every arc span, and $N_{v}$.

\paragraph{The estimator as a time average.}
With $\hat H=\hat H_{\mathrm{el}}+\hat H_{\mathrm{ph}}+\hat H_{\mathrm{int}}$, the key
move is \emph{not} to read the energy only at the endpoint $\tau=\beta$ but to use the
ground-state-dominated identity
\begin{equation}
  \langle\hat H\rangle \;=\; \frac1\beta\int_{0}^{\beta}\langle\hat H(\tau)\rangle\,d\tau,
  \label{eq:timeavg}
\end{equation}
i.e.\ the \emph{local} energy is averaged over the entire worldline. This is the
variance reduction: $\beta$ times more information per configuration than a single
asymptotic slope. Each term of $\hat H$ maps to one accumulator, and the three sum to
\begin{equation}
  E_{\mathrm{pol}}
  \;=\; \frac{1}{\beta\,\langle s\rangle}\,
  \Bigl\langle\;
    \underbrace{\sum_{\mathrm{seg}}(\varepsilon_{i}+\mu)\,\Delta\tau_{i}}_{\langle H_{\mathrm{el}}\rangle}
    \;+\;
    \underbrace{\sum_{\mathrm{arcs}}\tfrac12\,\omega_{a}\lvert\Delta\tau_{a}\rvert}_{\langle H_{\mathrm{ph}}\rangle}
    \;-\;
    \underbrace{(N_{v}-1)}_{\langle H_{\mathrm{int}}\rangle}
  \;\Bigr\rangle,
  \label{eq:estimator}
\end{equation}
with $\langle s\rangle$ the average sign. The three pieces are:
\begin{description}
  \item[Band energy $\langle H_{\mathrm{el}}\rangle$.] The code inserts the operator via
    \texttt{Hpropagator}, whose diagonal is $e^{-\varepsilon\tau}(\varepsilon+\mu)\tau$.
    The extra factor $\varepsilon\,\Delta\tau$ is exactly
    $\int_{0}^{\Delta\tau}\!ds\,\varepsilon\,e^{-\varepsilon\Delta\tau}
     =\varepsilon\,\Delta\tau\,e^{-\varepsilon\Delta\tau}$ --- the ``$\hat H e^{-\hat H\tau}$''
    insertion integrated over the segment. Summed over segments and divided by the
    diagram weight it yields $\tfrac1\beta\int_{0}^{\beta}\varepsilon_{k(\tau)}\,d\tau$.
    Energies are stored band-bottom subtracted ($\varepsilon$) with $\mu=\min\varepsilon$,
    so $\varepsilon+\mu$ restores the true band energy.
  \item[Phonon energy $\langle H_{\mathrm{ph}}\rangle$.] Each arc present over
    $\Delta\tau_{a}$ contributes $\omega_{a}$ to the instantaneous phonon energy;
    time-averaging gives $\tfrac1\beta\sum_{a}\omega_{a}\Delta\tau_{a}$. The $\tfrac12$
    corrects double counting of an internal arc seen from both its vertices; external
    (head/tail-anchored) phonon lines carry no $\tfrac12$.
  \item[Interaction energy $\langle H_{\mathrm{int}}\rangle$.] In the bare imaginary-time
    expansion the coupling-energy estimator evaluates to $-\langle N_{v}\rangle/\beta$
    (each vertex contributes $-1/\beta$); the ``$-1$'' offset is the head-vertex
    normalization convention.
\end{description}
Equation~\eqref{eq:estimator} is thus just
$\langle H_{\mathrm{el}}\rangle+\langle H_{\mathrm{ph}}\rangle+\langle H_{\mathrm{int}}\rangle
 =\langle\hat H\rangle$, measured at $\tau=\beta$ with no fitting window.

\paragraph{Finite-$\beta$ phonon propagator.}
The arcs use the thermal bosonic propagator
\begin{equation}
  D(\omega,\tau) \;=\; \frac{e^{-\omega\tau}+e^{-\omega(\beta-\tau)}}{1-e^{-\omega\beta}},
  \label{eq:dph}
\end{equation}
whose first term $e^{-\omega\tau}$ is spontaneous emission (surviving at $T=0$), second
term $e^{-\omega(\beta-\tau)}$ is thermal absorption / wrap-around, and denominator
$1/(1-e^{-\omega\beta})=1+n_{B}(\omega)$ is the Bose enhancement. As $\beta\to\infty$ it
collapses to $e^{-\omega\tau}$, which is what the code's \texttt{zeroTMC} flag
hard-codes --- and the LiF example uses \texttt{zeroTMC}, so LiF(e) is effectively the
$T=0$ limit even though the machinery is finite-$T$. Soft modes below a
$\texttt{phfreq\_cutoff}\approx1$~meV are dropped ($D\to0$) to avoid the
$1/(1-e^{-\omega\beta})$ and $(2M_{\kappa}\omega)^{-1/2}$ divergences.

\paragraph{Matrix (multiband) form.}

In multiband mode the segment propagators, vertices, and the $\hat H$-insertion are all matrices, so the diagram weight is $\operatorname{Tr}\!\bigl[\prod_{i}g_{i}\,e^{-H\Delta\tau_{i}}\bigr]$ and inserting $\hat H$ at one segment needs the product of everything to its left and right (the \texttt{left\_env}/\texttt{right\_env} environment matrices), with per-step rescaling. The sign is carried as $\operatorname{sign}(\operatorname{Re}\operatorname{Tr})$ and divided out via $\langle s\rangle$. In single band all of this collapses to the scalars of Eq.~\eqref{eq:estimator}; the trace is trivial.

\paragraph{Why it replaced the decay fit.}

Relative to the decay fit of P9, Eq.~\eqref{eq:estimator} removes the $\tau_{\max}$ window arbitrariness (the whole $[0,\beta]$ worldline is used), lowers variance (it time-averages rather than reading one slope), and has a systematic error controlled by $\beta$ (excited-state contamination $\sim e^{-\Delta\beta}$). The single-band form accumulates the three sums of Eq.~\eqref{eq:estimator} per sample, dividing by $\beta$ and $\langle s\rangle$ --- no environment matrices are needed until the multiband (hole/degenerate) work of D2. It has been cross-validated against the Fr\"ohlich $\alpha=1,2$ Mishchenko values (stage~0), where it agrees with the decay estimator, and against the exactly solvable dispersionless lattice model noted above; all production numbers (Section~\ref{sec:discretization}) use it.

\section{Phonon-cloud estimators from central time slices}
\label{sec:slice}

The energy estimator is exact at every $\tau$ because it is a log-derivative, $E_{\mathrm{loc}}=-\partial_{\tau}\ln G$. Occupation-type observables have no such identity. Inserting an operator $\hat O$ at time $t$ on the worldline gives
\begin{equation}
  O(t) = \frac{\langle0|c_{\mathbf p}\,e^{-\hat H(\tau-t)}\,\hat O\,e^{-\hat Ht}\,c^{\dagger}_{\mathbf p}|0\rangle}{G(\mathbf p,\tau)}
  \;\xrightarrow{\;t,\;\tau-t\,\gg\,1/\Delta\;}\;\langle\Psi_{0}(\mathbf p)|\hat O|\Psi_{0}(\mathbf p)\rangle ,
  \label{eq:slice}
\end{equation}
and the ground-state value is reached only far from both ends. In the bare-basis expansion the intermediate state at $t$ is the electron segment crossing $t$ plus the arcs spanning $t$. Hence, for $\hat O=\hat P_{N}$ (projector on $N$ phonons), $O(t)$ is the sign-weighted probability that exactly $N$ arcs cross $t$, i.e.\ FEP-DMC's $Z_{N,\mathbf p}$; $\hat N_{\mathrm{ph}}$ and its per-branch parts count the crossing arcs, and $\hat n_{\kv}$ reads the crossing segment's momentum ($f_{e}(\kv)$ of their Eq.~(4)). None of these needs external phonon lines or a new update.

The averaging window must be central. The whole-worldline average $\tau^{-1}\sum_{a}\Delta\tau_{a}$ is biased low by the ends, where the cloud is still building from $c^{\dagger}|0\rangle$. In the dispersionless model $N(t)=\bar N(1-e^{-\omega t})(1-e^{-\omega(\tau-t)})$, so the full average is $\bar N[1-2(1-e^{-\omega\tau})/(\omega\tau)+e^{-\omega\tau}]$: an $O(1/\omega\tau)$ error, about $-11\%$ for LiF at $\tau=240$~eV$^{-1}$. The production estimator averages Eq.~\eqref{eq:slice} over $t\in[\tau/2-w,\tau/2+w]$ (default $w=\tau/8$), which converges as $e^{-\Delta(\tau/2-w)}$ and is certified by $w$-independence. Implementation plan: \texttt{plans/open/nph-estimator.md}.

\section{Polaron dispersion, effective mass, and the acoustic-emission threshold}
\label{sec:dispersion}

\paragraph{Protocol.} $E_{0}(\mathbf p)$ is measured on grid momenta along $\Gamma\to X$ ($\mathbf p=\tfrac{j}{n}(\mathbf b_{1}+\mathbf b_{3})$) and $\Gamma\to L$, each point a separate fixed-$\tau$ run (\texttt{ironmc dispersion sweep}; runbook \texttt{notebooks/LiF-dispersion-mass.ipynb}). Cubic symmetry forces $E_{0}(\mathbf p)=E_{0}+|\mathbf p|^{2}/2m_{p}+c_{4,\mathrm{leg}}|\mathbf p|^{4}+\dots$ with one scalar $m_{p}$ (\texttt{latex/symmetry-paths}). Fit windows are fixed in physical $|\mathbf p|$ (\AA$^{-1}$), never in grid index: at fixed $j$ the momentum shrinks as $1/n$, and an index window mistakes non-parabolicity for grid dependence. The mass is reported as $m_{p}/m_{b}$, where the band mass $m_{b}=0.6926\,m_{e}$ is the $\Gamma$ curvature of the same Wannier $\varepsilon(\kv)$ the sampler uses (\texttt{examples/band\_mass.rs}). FEP-DMC's ``$m_{p}/m_{e}$'' uses exactly this normalization; $m_{e}$ there is the band effective mass.

\paragraph{Current result.} With the $p^{2}+p^{4}$ model over $|\mathbf p|\le0.16$~\AA$^{-1}$: $m_{p}/m_{b}=5.31\pm0.23\,(\text{stat})\pm0.42\,(\text{window})$ at $80^{3}$ and $5.0\pm0.3\pm1.0$ at $120^{3}$. At the same $|\mathbf p|$, $40^{3}$, $80^{3}$ and $120^{3}$ agree to within about 1~meV, so the dispersion converges in $n$ much faster than $E_{0}$. A pure parabola is excluded beyond $|\mathbf p|\approx0.1$~\AA$^{-1}$ ($\chi^{2}$ grows, the fitted mass drifts upward): parabolic fits give $6.0$--$6.4$ over $|\mathbf p|\le0.12$~\AA$^{-1}$ and up to $\approx10$ over wider ranges. FEP-DMC's 6.1 ($\Gamma$--$L$, parabolic fit) is therefore consistent with our dispersion, but a single parabolic number is window-dependent at the $\pm1$ level.

\paragraph{The acoustic-emission (Landau) threshold.} For each total momentum $\mathbf p$, the fixed-$\tau$ estimator converges as $\tau\to\infty$ to the \emph{lowest} energy in that sector:
\begin{equation}
  E_{\mathrm{sector}}(\mathbf p)=\min\Bigl\{E_{\mathrm{qp}}(\mathbf p),\;\min_{\qv,\nu}\bigl[E_{\mathrm{qp}}(\mathbf p-\qv)+\omega_{\nu}(\qv)\bigr],\;\dots\Bigr\}.
\end{equation}
The second entry is the polaron recoiling to $\mathbf p-\qv$ after emitting one real phonon. For LO phonons it is the familiar flattening at $E_{0}+\omega_{\mathrm{LO}}$. For acoustic phonons, $\omega\approx c_{s}q$ with no gap. With $E_{\mathrm{qp}}\approx E_{0}+p^{2}/2m_{p}$, the emission state lies below the quasiparticle as soon as the polaron's group velocity exceeds the sound velocity, $|\mathbf p|/m_{p}>c_{s}$: the lattice analogue of Cherenkov radiation. Beyond that momentum the bottom of the sector is $\approx E_{0}+c_{s}|\mathbf p|-m_{p}c_{s}^{2}/2$. For LiF, $\omega_{\mathrm{TA}}=4.93$~meV at $|\qv|=0.155$~\AA$^{-1}$ along $\Gamma\to X$ gives $c_{s}\approx32$~meV\,\AA{} ($4.8$~km/s). With $m_{p}\approx3.7\,m_{e}$ this puts the threshold at $p_{c}=m_{p}c_{s}\approx0.015$~\AA$^{-1}$, below the first nonzero $\Gamma\to X$ grid point even at $120^{3}$ ($0.026$~\AA$^{-1}$).

The quasiparticle becomes a resonance that can decay by acoustic emission. The measured $E_{0}(\mathbf p,\tau)$ remains quasiparticle-dominated only while the admixture of the lower continuum, $\sum_{i}w_{i}e^{+\delta_{i}\tau}$ with $\delta_{i}=E_{\mathrm{qp}}-E_{i}>0$, stays small. Each weight is small, $w_{i}\sim|g_{\mathrm{ac}}|^{2}/(N_{q}\delta_{i}^{2})$ (acoustic coupling is weak and each state carries $1/N_{q}$), but the factor $e^{\delta_{i}\tau}$ grows with $\tau$ and with $|\mathbf p|$. So, unlike at $\Gamma$, longer $\tau$ is not automatically better at finite $\mathbf p$. The diagnostic is cheap: at the same $(n,\mathbf p)$, $E_{0}(\mathbf p,\tau240)-E_{0}(\mathbf p,\tau60)$ must equal the $\Gamma$ shift ($-2.8$~meV) if the quasiparticle dominates, while an excess that grows with $|\mathbf p|$ signals contamination. Keeping mass fits at small $|\mathbf p|$, where $\delta$ is a few meV, limits the effect in any case.

\section{Remote bands: the Allen--Heine--Cardona correction}
\label{sec:ahc}

The DiagMC is exact for the Hamiltonian $\hat H_{n}$ it samples: one Wannier band (the CBM), linear e--ph coupling. Two pieces of the physical band-edge renormalization are outside it: (i) Fan--Migdal virtual transitions into bands not included in the sampled manifold, and (ii) the Debye--Waller term (second order in the displacement at one vertex), which a linear-coupling DiagMC never generates. FEP-DMC adds both perturbatively (AHC, their ref.~36, 37), moving LiF(e) from $-0.408$ to $-0.538$~eV. This is a 130~meV correction, larger than every DiagMC systematic above. Consistently with the sampled manifold,
\begin{equation}
  \Delta E_{\mathrm{AHC}} = \mathrm{Re}\,\Sigma^{\mathrm{FM}}_{\mathrm{CBM}}\bigl(\varepsilon_{\mathrm{CBM}}\bigr)\Big|_{m\notin\{\mathrm{CBM}\}} + \Sigma^{\mathrm{DW}}_{\mathrm{CBM}}, \qquad
  \Sigma^{\mathrm{FM}}_{n\kv}(\varepsilon)=\frac{1}{N_{q}}\sum_{m\nu\qv}|g_{mn\nu}(\kv,\qv)|^{2}\,\frac{1}{\varepsilon-\varepsilon_{m\kv+\qv}-\omega_{\nu\qv}} ,
\end{equation}
at $T=0$ (single carrier). The band sum runs over all Kohn--Sham bands, via Sternheimer for the high ones, and excludes the intraband CBM$\to$CBM term that the DiagMC already contains to all orders. The practical route and its pitfalls are in \texttt{plans/open/ahc-correction.md}.

\end{document}
